% answers_short.tex —— 2011 年·福建卷·文科 · 简短答案速查
% 由 tools/make_short_answers.py 自 answers.tex 的 \gksolhead/\gkanswer 提取，
% 并与 2011/2011_数学真题_福建文_答案要点.md 逐题比对后生成，勿手改（改源 MD 后重跑）。

\begin{gknotice}
\textbf{使用说明}：本节只给最终答案便于快速核对，完整推导见解析版。答案为编写者独立推导，\textbf{非官方标准答案}。
\end{gknotice}

\gkpart{选择题}
\begin{center}\begin{tabular}{ccccccccccccc}
\toprule
题号 & 1 & 2 & 3 & 4 & 5 & 6 & 7 & 8 & 9 & 10 & 11 & 12 \\
\midrule
答案 & A & D & A & B & B & C & C & A & D & D & A & C \\
\bottomrule
\end{tabular}\end{center}

\gkpart{填空题}
\begin{itemize}[leftmargin=2.6em]
  \item[\textbf{13.}] $1$。
  \item[\textbf{14.}] $2$。
  \item[\textbf{15.}] $\sqrt2$。
  \item[\textbf{16.}] $\dfrac{\sqrt5-1}2$。
\end{itemize}

\gkpart{解答题}
\begin{itemize}[leftmargin=2.6em]
  \item[\textbf{17.}] (1) $a_n=3-2n$；(2) $k=7$
  \item[\textbf{18.}] (1) $b=-1$；(2) $(x-2)^2+(y-1)^2=4$
  \item[\textbf{19.}] (1) $a=0.1$，$b=0.15$，$c=0.1$；(2) 所有结果见解析；概率 $\dfrac25$
  \item[\textbf{20.}] (1) 证明见解析；(2) $\dfrac{\sqrt3-\sqrt2+\sqrt6}6$
  \item[\textbf{21.}] (1) $2$；(2) $\theta\in\left[0,\dfrac\pi2\right]$，最小值 $1$，最大值 $2$
  \item[\textbf{22.}] (1) $b=2$；(2) $a>0$ 时减区间 $(0,1)$、增区间 $(1,+\infty)$；$a<0$ 时增区间 $(0,1)$、减区间 $(1,+\infty)$；(3) 存在，$m=1$，$M=2$
\end{itemize}
